\section{System Architecture}

A \(512\times512\) grayscale cover carries a \(128\times128\) grayscale secret. Inputs enter the digital path as uint8 images, and clipping plus half-up rounding define the output boundary. The four most-significant bits of each secret pixel form Base; the four least-significant bits form Detail. Each layer contains 8,192 raw bytes.

Figure~\ref{fig:pipeline} shows the complete path. The cover determines transform coordinates, the secret determines the layered payload, and a calibration profile contributes band-stability evidence. These three streams meet only at the embedding plan. At the receiver, the original cover is used for semi-blind coefficient differencing, after which Reed--Solomon and CRC checks decide whether a secret image may be emitted.

\begin{figure*}[t]
\centering
\resizebox{\textwidth}{!}{%
\begin{tikzpicture}[
  font=\small,
  box/.style={draw, rounded corners=2pt, align=center, minimum height=9mm, minimum width=25mm, fill=palegray},
  cover/.style={box, fill=baseblue!10},
  secret/.style={box, fill=accentorange!12},
  calib/.style={box, fill=softgreen!12},
  decision/.style={box, fill=softgreen!10, minimum width=29mm},
  arrow/.style={-{Latex[length=2.2mm]}, thick},
  dashedarrow/.style={-{Latex[length=2.2mm]}, thick, dashed}
]
\node[cover] (cover) at (0,2.4) {Cover image};
\node[cover] (pdfb) at (3.2,2.4) {PDFB analysis};
\node[cover] (coords) at (6.5,2.4) {Eligible\\coordinates};

\node[secret] (secret) at (0,0.8) {Secret image};
\node[secret] (split) at (3.2,0.8) {Base/Detail\\split};
\node[secret] (ecc) at (6.5,0.8) {RS coding, scrambling\\and interleaving};

\node[calib] (calib) at (0,-0.8) {Calibration covers};
\node[calib] (stability) at (3.2,-0.8) {Band stability};
\node[calib] (score) at (6.5,-0.8) {Adaptive score\\and quota};

\node[box, fill=baseblue!8, minimum width=28mm] (embed) at (10.0,1.1) {Deterministic\\embedding plan};
\node[box] (stego) at (13.0,1.1) {Stego image};
\node[box, fill=softred!8] (channel) at (15.8,1.1) {Digital channel};
\node[box] (extract) at (18.8,1.1) {Semi-blind\\extraction};
\node[decision] (validate) at (22.0,1.1) {Header, RS and\\CRC validation};
\node[secret, minimum width=27mm] (output) at (25.3,1.1) {Recovered secret\\or no output};

\draw[arrow] (cover) -- (pdfb);
\draw[arrow] (pdfb) -- (coords);
\draw[arrow] (secret) -- (split);
\draw[arrow] (split) -- (ecc);
\draw[arrow] (calib) -- (stability);
\draw[arrow] (stability) -- (score);
\draw[arrow] (coords) -- (embed);
\draw[arrow] (ecc) -- (embed);
\draw[arrow] (score) -- (embed);
\draw[arrow] (embed) -- (stego);
\draw[arrow] (stego) -- (channel);
\draw[arrow] (channel) -- (extract);
\draw[arrow] (extract) -- (validate);
\draw[arrow] (validate) -- (output);
\draw[dashedarrow] (cover.south) .. controls (8,-1.8) and (16,-1.8) .. (extract.south);
\node[below=1mm of extract, font=\scriptsize, text=black!65] {original cover};
\end{tikzpicture}%
}
\caption{Audited hierarchical transport. Base and Detail are coded independently; calibration affects coordinate ranking but does not use the final test pairs. A recovered image is emitted only after the complete integrity contract succeeds.}
\label{fig:pipeline}
\end{figure*}

\subsection{Hierarchical unequal protection}

C0 and C1 protect Base and Detail symmetrically. Each layer receives 32 RS(255,127) and 22 RS(255,191) codewords. C2 and C3 keep the same total number of codewords but redistribute them: Base receives 65 RS(255,127) blocks and Detail receives 43 RS(255,191) blocks. The design does not create extra channel capacity; it spends more of the existing redundancy on the bits that dominate the reconstructed pixel value.

The nominal Base correction budgets are
\begin{equation}
\begin{aligned}
C0 &: \quad 32(64)+22(32)=2,752,\\
C3 &: \quad 65(64)=4,160
\end{aligned}
\qquad \text{symbols}.
\end{equation}
The increase is 51.2\%, subject to the usual per-codeword condition that no block exceeds its own correction radius. Unit tests correct exactly 64 versus fail at 65 altered symbols for RS(255,127), and correct 32 versus fail at 33 for RS(255,191).

\subsection{Transform-bound adaptive allocation}

Adaptive methods score each eligible band from coefficient energy, variance, entropy of absolute magnitudes, and channel stability measured on designated calibration covers. Median/MAD normalisation limits the effect of outlying bands. The stability artifact records both the calibration attacks and the transform fingerprint, and the final runner refuses to apply it to a different transform.

Scores control integer band quotas and local embedding weights between 0.75 and 1.25. C1 changes allocation but retains symmetric coding. C2 changes coding but retains fixed allocation. C3 combines both and writes Base before Detail in descending score order.

\subsection{Capacity and quality controls}

The allocator selects exactly 222,360 unique coefficient slots and cannot exceed the capacity of any eligible band. A binary search then finds the largest global embedding strength that remains feasible after inverse transformation, clipping, rounding, and uint8 conversion. All variants therefore carry the same payload and meet the same realised cover--stego PSNR target of \SI{45.0\pm0.1}{dB}. Supplementary Notes S1--S3 provide the header layout, seed derivation, allocation invariants, and coefficient-map audit.
